%% ==========================================================
%% Springer Nature Journal Article — sn-jnl class
%% Math and Physical Sciences Numbered Reference Style
%% ==========================================================
%% Changes from original:
%%   1. Citations [1]–[10] added at end of relevant paragraphs
%%   2. Bibliography author names corrected per actual PDFs
%%   3. Algorithm Used column added to literature survey table
%%   4. All user-written content preserved exactly
%% ==========================================================

\documentclass[pdflatex,sn-mathphys-num]{sn-jnl}

%%%% Standard Packages %%%%
\usepackage{graphicx}
\usepackage{multirow}
\usepackage{amsmath,amssymb,amsfonts}
\usepackage{amsthm}
\usepackage{mathrsfs}
\usepackage[title]{appendix}
\usepackage{xcolor}
\usepackage{textcomp}
\usepackage{manyfoot}
\usepackage{booktabs}
\usepackage{algorithm}
\usepackage{algorithmicx}
\usepackage{algpseudocode}
\usepackage{listings}
\usepackage{tabularx}
\usepackage{threeparttable}
\usepackage[T1]{fontenc}
\usepackage[utf8]{inputenc}
\usepackage{lmodern}

\theoremstyle{thmstyleone}
\newtheorem{theorem}{Theorem}
\newtheorem{proposition}[theorem]{Proposition}

\theoremstyle{thmstyletwo}
\newtheorem{example}{Example}
\newtheorem{remark}{Remark}

\theoremstyle{thmstylethree}
\newtheorem{definition}{Definition}

\raggedbottom

%% ==========================================================
%% DOCUMENT START
%% ==========================================================

\begin{document}

\title[SmartSpace: Warehouse Recommendation]{Intelligent Warehouse Recommendation
for Maharashtra Using Hybrid Machine Learning
and LLM Assistance}

%%% Authors %%%

\author*[1]{\fnm{Varsha} \sur{Wangikar}}
\email{ksdeherkar@apsit.edu.in}

\author[1]{\fnm{Dipesh} \sur{Sharma}}
\email{dipeshsharma34@apsit.edu.in}

\author[1]{\fnm{Pranjal} \sur{Zambre}}
\email{pranjalzambre115@apsit.edu.in}

\author[1]{\fnm{Manthan} \sur{Shinde}}
\email{manthanshinde120@apsit.edu.in}

\author[1]{\fnm{Sahil} \sur{Shinde}}
\email{sahilshinde13@apsit.edu.in}

%%% Affiliation %%%

\affil*[1]{%
  \orgdiv{Department of Computer Engineering},
  \orgname{A.P. Shah Institute of Technology},
  \orgaddress{%
    \city{Thane},
    \country{India}%
  }%
}

%%% Abstract %%%

\abstract{
We built SmartSpace, a platform that helps people find and pick
the right warehouse in Maharashtra, India.
It works with a database of more than 10,000 warehouse records.
The core of the system is a hybrid machine learning setup
where K-Nearest Neighbors and Random Forest models feed into
a meta-learner, and the output is a stable, reliable warehouse ranking.
A weighted scoring method looks at things like how close the warehouse is,
what it costs, how much space it has, what type of facility it is,
whether it is available, and how users have rated it.
We also brought in Meta Llama 3.3 70B as a large language model
so the platform can give clear explanations and suggest alternatives
in plain language.
If the LLM becomes unavailable for any reason, the system falls back
to local ML models and keeps working without interruption.
The whole thing runs on Supabase PostgreSQL, React TypeScript,
and Express.js.
In our early tests, SmartSpace cut user search time by about 60\%
and hit a recommendation success rate above 85\%.
We think this kind of setup, where ML handles the numbers and an LLM
handles the reasoning, is a practical way forward for warehouse
selection tools.
}

%%% Keywords %%%

\keywords{
warehouse recommendation,
hybrid machine learning,
random forest,
K-nearest neighbors,
large language models,
intelligent logistics,
supply chain optimization
}

\maketitle

%% ==========================================================
%% SECTION 1 — INTRODUCTION
%% ==========================================================

\section{Introduction}\label{sec1}

Maharashtra has witnessed significant growth in manufacturing,
retail distribution, and e-commerce operations over the past decade,
leading to increased demand for flexible and location-specific
warehousing solutions.
Despite this growth, the process of identifying suitable warehouses
remains largely manual, requiring businesses to consult brokers,
browse fragmented registries, and conduct repeated site visits.
This approach frequently results in extended search durations,
mismatched facility selections, and elevated operational costs,
particularly for small and medium enterprises
\cite{b1, b5, b8}.

The challenge of warehouse selection extends beyond availability
alone.
Decision-makers must simultaneously consider geographic proximity,
pricing constraints, storage capacity, facility specialization,
compliance status, and service quality.
Existing digital platforms typically address only isolated aspects
of this decision-making process and lack the capability to provide
multi-criteria recommendations with transparent explanations
\cite{b2}.
Additionally, newly listed or infrequently booked warehouses often
suffer from limited historical data, making accurate recommendations
difficult through traditional collaborative filtering approaches
\cite{b10}.

SmartSpace addresses these limitations through an intelligent
warehouse recommendation platform that integrates structured
machine learning techniques with large language model assistance.
The system matches warehouse seekers with suitable facilities
by jointly analyzing explicit user preferences and implicit
constraints derived from conversational input.
Machine learning models generate stable numerical rankings,
while the language model provides contextual explanations
and alternative suggestions to enhance user understanding
and trust \cite{b3, b7}.

The primary contributions of this work are fourfold.
First, a domain-specific scoring function is proposed that
consolidates location relevance, pricing compatibility,
storage adequacy, facility type, availability, verification status,
and quality ratings into a single interpretable metric
\cite{b1, b5}.
Second, a hybrid recommendation framework is developed by
combining K-Nearest Neighbors and Random Forest algorithms
through a meta-learning approach to improve ranking consistency
across diverse user profiles \cite{b10}.
Third, an LLM-assisted explanation layer is integrated to
support natural language interaction and semantic refinement
of recommendations, accompanied by a robust fallback mechanism
to ensure uninterrupted service \cite{b7, b4}.
Finally, a production-ready implementation is delivered using
Supabase PostgreSQL, React TypeScript, and a modular API-driven
backend, demonstrating practical applicability within real-world
logistics workflows \cite{b6}.

%% ==========================================================
%% SECTION 2 — LITERATURE SURVEY
%% ==========================================================

\section{Literature Survey}\label{sec2}

A fair amount of research has gone into warehouse management,
logistics platforms, and the use of AI in supply chain work.
Some of it looks at warehouse sharing models, others focus
on on-demand logistics, and a few recent papers explore how
large language models can help with planning and decision-making
in this space \cite{b7, b8}.

That said, most of this work tackles one piece of the puzzle
at a time.
One paper might deal with pricing, another with location
selection, and yet another with automation \cite{b2}.
Very few have tried to bring all of these together into a single
recommendation framework that also explains its reasoning
and has a backup plan when something goes wrong \cite{b3, b4}.
On top of that, many existing platforms struggle with what is
called the ``cold-start'' problem, where a newly listed warehouse
has no booking history, so the system does not know what to do
with it \cite{b10}.
We address this through content-based filtering that relies
on warehouse features rather than past user behavior.

Table~\ref{tab:literature_survey} lists the key studies
we reviewed, what they contributed, the algorithm each paper
actually uses, and where they fall short.
This gap in the literature, the lack of a system that combines
hybrid ML with explainable LLM reasoning and reliable fallback,
is what motivated our work.

%%% Literature Survey Table — Algorithm column updated %%%

\begin{table}[h]
\label{tab:literature_survey}
\centering
\footnotesize
\renewcommand{\arraystretch}{1.25}
\setlength{\tabcolsep}{2pt}
\begin{tabular}{|c|p{4cm}|p{3cm}|p{4cm}|}
\hline
\textbf{Paper Title}
  & \textbf{Author \& Publication}
  & \textbf{Algorithm Used}
  & \textbf{Findings / Limitations} \\
\hline

Enhanced Modeling for Warehouse Sharing Platform [1]
  & Z. Xing et al., arXiv (2024)
  & Adaptive Large Neighbourhood Search (ALNS) Heuristic
  & Framework for warehouse-sharing platforms;
    high latency for real-time inference on large datasets. \\
\hline

Shared Logistics, Literature Review [2]
  & M. Matusiewicz, D. Ksi\k{a}\.zkiewicz,
    Applied Sciences (2023)
  & Systematic Literature Review
  & Reviews trends and challenges in coordination,
    trust, and data transparency;
    limited to structured data. \\
\hline

LLMs for Simulation-Based Optimization in SCM [3]
  & T. Wi\'{s}niewski, APEM (2025)
  & LLM-assisted Simulation Optimization (ChatGPT)
  & LLMs enhance planning and service;
    prone to hallucination,
    lacks domain-specific grounding. \\
\hline

Decentralized Warehouse Management Using LLMs [4]
  & T. Berlec et al.,
    Applied Sciences (2025)
  & LLM-driven Multi-Agent System
    (Game Theory, Decision Theory)
  & Decentralized LLM-based warehouse agents
    outperform random strategies;
    conceptual, no integrated recommendation. \\
\hline

On-Demand Warehousing Platforms [5]
  & V. Elia et al., Logistics (2024)
  & Trend Analysis (Systematic Review)
  & Impact on scalability and operational
    efficiency analyzed;
    cannot model non-linear pricing. \\
\hline

Full Stack Systems Report [6]
  & Technical Report,
    React and Supabase
  & Full Stack Development
    (React, Supabase, Express.js)
  & Implementation reference for web-based
    logistics platforms. \\
\hline

LLMs for Supply Chain Management [7]
  & H. Wang et al., arXiv (2025)
  & RAG Framework,
    Domain-specialized LLM,
    Game-theoretic Analysis
  & Improved forecasting and decisions
    via LLM analytics;
    dependency on external APIs,
    latency concerns. \\
\hline

Warehousing 5.0 [8]
  & B. Y. Ekren et al., IJPR (2026)
  & Conceptual Framework
    (Warehousing 5.0 Vision)
  & Future vision for intelligent and adaptive
    warehouse ecosystems;
    conceptual, no empirical validation. \\
\hline

Smart Logistics Warehouse Mgmt.\ System [9]
  & C. Liang et al., ICJE (2025)
  & Multi-modal AI
    (DeepSeek-R1 + GLM-4V)
  & Multi-modal AI for smart logistics;
    complex orchestration,
    system instability during failovers. \\
\hline

ML in Warehouse Mgmt.\ Survey [10]
  & R. F. de Assis et al.,
    Procedia CS (2024)
  & Systematic Literature Review
    (ML Methods Survey)
  & ML for demand forecasting, inventory control,
    space optimization;
    gap in combined ML + LLM systems. \\
\hline
\end{tabular}
\begin{center}
    \textbf{Table 1} Literature Survey
\end{center}
\end{table}

%% ==========================================================
%% SECTION 3 — METHODOLOGY
%% ==========================================================

\section{Methodology}\label{sec3}

SmartSpace was designed in layers.
Each layer does its own job, but they all feed into one
decision-support system.
We kept structured data processing separate from
language-based reasoning on purpose, so that if one part
needs to change, it does not break everything else.
The system can scale because of this separation,
and it handles different kinds of data without mixing
things up \cite{b4, b9}.

Here is how the pipeline works, step by step:
\begin{itemize}
    \item \textbf{Data Acquisition:}
      We collected government-certified warehouse registry
      data from Indian regulatory authorities
      \cite{b1, b5}.

    \item \textbf{Data Preprocessing:}
      The raw data went through cleaning, geospatial
      gap-filling, and synthetic augmentation so it
      looks like what you would actually see in the market
      \cite{b10}.

    \item \textbf{Model Selection and Training:}
      We implemented a 5-algorithm ML ensemble for
      numerical scoring and plugged in a Large Language
      Model for reasoning over text
      \cite{b3, b7}.

    \item \textbf{Multimodal Fusion:}
      We combined what users explicitly ask for with what
      they seem to want based on their conversation, and
      used both signals to rank warehouses
      \cite{b9}.

    \item On the architecture side, there is a React-based
      front end for the user interface.
      Behind it sits an application layer that handles
      API routing and authentication.
      The real work happens in what we call the
      Intelligence Layer, which has two engines:
      Meta Llama 3.3 70B for reasoning tasks and a
      five-model ML ensemble (KNN, RF, GB, NN, and
      Stacking) for number-based predictions \cite{b4, b10}.
      Everything is stored in Supabase PostgreSQL with
      row-level security turned on, and we also pull in
      external WDRA regulatory data \cite{b6}.
\end{itemize}

\begin{figure}[h]
\centering
\includegraphics[width=0.75\columnwidth]{architecture.jpeg}
\begin{center}
    \textbf{Fig. 1} System Architecture of SmartSpace
\end{center}
\label{fig:arch}
\end{figure}

%%% 3.1 Data Acquisition %%%

\subsection{Data Acquisition}\label{subsec:data-acq}

We started with data from the
\textit{Warehousing Development and Regulatory Authority (WDRA)},
which is a government body under India's Department of Food
and Public Distribution.
This gave us credibility right from the start \cite{b5}.
\begin{itemize}
    \item \textbf{Structured Registry Data:}
      The base dataset has 424 certified warehouse records
      spread across 33 districts in Maharashtra.
      Each record includes the warehouse capacity in metric
      tons, its district, registration validity period,
      and contact details.

    \item \textbf{Market Context:}
      Government data alone was not enough.
      We added market-specific features like pricing
      (INR per square foot per month) and facility types
      such as Cold Chain, FMCG, and Industrial.
      These came from regional logistics reports so the
      data would reflect actual market conditions
      \cite{b1}.
\end{itemize}

%%% 3.2 Data Preprocessing %%%

\subsection{Data Preprocessing}\label{subsec:preproc}

The raw data from WDRA needed a lot of work before we
could use it for machine learning \cite{b10}.
\begin{itemize}
    \item \textbf{Geospatial Imputation:}
      The original dataset did not have precise latitude
      and longitude values.
      We used a centroid-based method, assigning each
      warehouse the coordinates of its district center
      and then adding a small random Gaussian offset
      ($\pm 0.05^\circ$) to spread them out realistically.

    \item \textbf{Data Cleaning:}
      We used regex-based checks to clean up phone numbers
      and pincodes.
      If a warehouse had an expired license, we flagged it
      and lowered its health score so the recommendation
      engine would rank it lower.
\end{itemize}

%%% 3.3 Dataset Description and Preparation %%%

\subsection{Dataset Description and Preparation}%
\label{subsec:dataset}

424 records is not enough to train a recommendation engine
that can handle diverse user preferences.
So we augmented the data programmatically to reach
\textbf{10,002 records} \cite{b10}.
\begin{itemize}
    \item \textbf{Feature Engineering:}
      We grew the feature set from 12 attributes to 16.
      The new ones include occupancy rates, quality ratings
      (a composite of amenities and certifications), and
      whether the warehouse offers micro-rental spaces.

    \item \textbf{Normalization:}
      Features like \texttt{total\_area} and
      \texttt{price\_per\_sqft} were scaled using Min-Max
      normalization.
      This was needed so that the KNN algorithm would not
      give too much weight to features with large numeric
      ranges.

    \item \textbf{Class Balancing:}
      We used stratified sampling to make sure all
      7 warehouse types (Pharma, Agriculture, General,
      and so on) were represented fairly.
      Without this step, the model would have been biased
      toward the most common types.
\end{itemize}

%%% 3.4 Model Selection and Training %%%

\subsection{Model Selection and Training}\label{subsec:model}

We went with a dual-engine design.
One engine handles the numbers, the other handles language
\cite{b3, b10}.

\subsubsection{Structured Recommendation Engine (Ensemble)}%
\label{subsubsec:ensemble}

For matching warehouses to user preferences based on numbers
for things like price, area, and location, we built a
\textbf{5-algorithm ensemble}.
Each algorithm scores every warehouse independently, and then
a weighted vote decides the final ranking \cite{b10}.
\begin{itemize}
    \item \textbf{K-Nearest Neighbors (KNN):}
      We picked KNN because it works well for spatial
      clustering problems.
      It measures how far each warehouse is from the user's
      ideal warehouse using a distance formula in feature
      space.

    \item \textbf{Random Forest:}
      This was chosen because it handles mixed data types
      well, especially categorical ones like
      \texttt{warehouse\_type} and \texttt{amenities}.
      When we looked at feature importance, location and
      price together accounted for about 75\% of the
      decision weight.
\end{itemize}

\subsubsection{Semantic Reasoning Engine (LLM)}%
\label{subsubsec:llm}

For understanding what users actually mean when they type
in natural language, we use \textbf{Meta Llama 3.3 70B}
as our language model \cite{b3, b7, b8}.
\begin{itemize}
    \item \textbf{Intent Parsing:}
      The LLM pulls out key pieces from what the user types.
      If someone writes ``Find cold storage in Pune under
      50k,'' the model extracts the location (Pune),
      budget (50k), and constraint (cold storage) \cite{b7}.

    \item \textbf{Re-Ranking:}
      After the ML ensemble returns its top 50 matches,
      the LLM looks at them again.
      It re-ranks them based on things the user might care
      about but did not explicitly say \cite{b4}.
\end{itemize}

%%% 3.5 Multimodal Fusion %%%

\subsection{Multimodal Fusion}\label{subsec:fusion}

The ML numbers and the LLM reasoning come together in
what we call the \textbf{Interactive Preference Engine}.
Users set their hard preferences, like price range and area,
through form inputs.
At the same time, they can add softer preferences like
``Prefer verified facilities'' through conversation.
These soft preferences adjust the internal weights dynamically.
The result is a recommendation list that is both
mathematically sound and contextually relevant
\cite{b9, b4}.

%% ==========================================================
%% SECTION 4 — IMPLEMENTATION
%% ==========================================================

\section{Implementation}\label{sec4}

SmartSpace is a working system, not a prototype on paper.
It runs as a full-stack web application that handles three
kinds of users: Seekers who look for warehouses, Owners
who list them, and Admins who oversee the whole process.
Everything was built to respond fast enough for real use
\cite{b6}.

%%% 4.1 Technology Stack %%%

\subsection{Architectural Framework and Technology Stack}%
\label{subsec:tech}

We split the front end and back end so they can be worked
on independently.
\begin{itemize}
    \item \textbf{Frontend Layer:}
      The interface is built with \textbf{React 18.3} and
      \textbf{TypeScript} for type safety.
      We used \textbf{TanStack Query} to manage server state,
      which handles caching and keeps the UI in sync with the
      database.
      Styling is done with \textbf{TailwindCSS}, and the whole
      thing is bundled through \textbf{Vite 6.3}, which gives
      us fast page reloads during development.

    \item \textbf{Backend and Database Layer:}
      On the server side, we have \textbf{Node.js} running
      \textbf{Express.js 4.x} for API routing.
      The database is \textbf{Supabase PostgreSQL}, and we
      turned on \textbf{Row-Level Security (RLS)} so each user
      only sees what they are supposed to see.
      For real-time updates on bookings and inventory, we use
      Supabase's built-in WebSocket subscriptions \cite{b6}.
\end{itemize}

%%% 4.2 Functional Modules %%%

\subsection{Functional Module Implementation}%
\label{subsec:modules}

The application has three main modules, each aimed at
a different type of user.

\subsubsection{Seeker Module (Discovery and Transaction)}%
\label{subsubsec:seeker}

This is what most users interact with.
It has two AI subsystems plus a manual search option:
\begin{itemize}
    \item \textbf{Interactive ML Recommendation Engine:}
      The user fills in what they want, things like target
      price and minimum area, and the system builds a query
      vector for the KNN algorithm from those inputs
      \cite{b10}.

    \item \textbf{Smart Booking Assistant:}
      This uses a four-step NLP pipeline.
      The LLM reads unstructured text like ``I need 400 sq ft
      in Thane,'' pulls out the relevant details, and turns
      them into a search query the system can act on
      \cite{b3, b7}.

    \item \textbf{Parametric Discovery Engine:}
      For users who prefer structured search, this mode runs
      SQL queries directly against warehouse attributes like
      type, location, and amenities.
\end{itemize}

\subsubsection{Owner Module (Asset Management)}%
\label{subsubsec:owner}

Warehouse owners use this module to manage their listings:
\begin{itemize}
    \item \textbf{Dynamic Pricing Inference:}
      A pre-trained regression model suggests a
      price-per-square-foot range based on current market
      data from MIDC regions.
      Owners can accept or adjust \cite{b5}.

    \item \textbf{Inventory Dashboard:}
      Built with \textbf{Recharts}, it shows occupancy rates
      and revenue numbers.
      Notifications come in real time when bookings change.
\end{itemize}

\subsubsection{Admin Module (Governance)}%
\label{subsubsec:admin}

Admins approve or reject warehouse submissions and manage
user roles.
Role-Based Access Control (RBAC) keeps Owners and Seekers
separated in terms of what they can do on the platform
\cite{b6}.

%%% 4.3 Hybrid AI Engine %%%

\subsection{Hybrid AI Engine: ML Ensemble}%
\label{subsec:hybrid}

This is where the real intelligence sits.
Two engines work in parallel: a 5-model ML ensemble for
scoring and a language model for reasoning
\cite{b10, b3}.

\subsubsection{The 5-Algorithm ML Ensemble}%
\label{subsubsec:ml5}

All five algorithms run at the same time, and each one
produces a relevance score $S_i \in [0, 1]$ for every
warehouse-query pair.
The input is an 8-dimensional feature vector: location match,
price fitness, area adequacy, warehouse type, verification
status, availability, user rating, and amenity density.

\textbf{K-Nearest Neighbors (K=8).}
KNN computes a weighted Euclidean distance between the
warehouse feature vector
$\mathbf{x} = (x_1, x_2, \ldots, x_n)$
and the user's ideal vector
$\mathbf{x}^* = (x_1^*, x_2^*, \ldots, x_n^*)$:

\begin{equation}
d(\mathbf{x}, \mathbf{x}^*)
  = \frac{%
      \sqrt{\sum_{i=1}^{n} w_i \left(x_i - x_i^*\right)^2}%
    }{%
      \sqrt{\sum_{i=1}^{n} w_i}%
    }
\label{eq:knn}
\end{equation}

\noindent
Here, $w_i$ is the weight for each feature
(location: 0.30, price: 0.25, area: 0.20, type: 0.10,
verification: 0.05, availability: 0.05, rating: 0.03,
amenities: 0.02), and $n = 8$.
The similarity score is just
$S_{\text{KNN}} = 1 - d(\mathbf{x}, \mathbf{x}^*)$.

\textbf{Random Forest (50 Trees).}
Each of the 50 decision trees trains on a bootstrap sample
covering 80\% of the candidates, and each split picks from
$\lfloor\sqrt{n}\rfloor$ randomly chosen features.

\textbf{Gradient Boosting (10 Iterations).}
Trees are added one after another, each one correcting what
the previous ones got wrong.
The learning rate is $\eta = 0.1$:
$F_m(\mathbf{x}) = F_{m-1}(\mathbf{x})
  + \eta \cdot h_m(\mathbf{x})$.

\textbf{Neural Network
(10$\rightarrow$10$\rightarrow$4$\rightarrow$1).}
A four-layer feedforward network.
The hidden layers use LeakyReLU ($\alpha = 0.01$) and the
output layer uses sigmoid to squeeze the result between 0
and 1.

\textbf{Hybrid Stacking.}
A meta-learner that takes the KNN and Random Forest outputs
and combines them for a more stable prediction.

The ensemble score is a weighted average of all five:
\begin{equation}
E(w) = \sum_{j=1}^{5} \alpha_j \cdot S_j(w)
\label{eq:ensemble}
\end{equation}

\noindent
The weights are
$\alpha_{\text{KNN}} = 0.25$,
$\alpha_{\text{RF}} = 0.25$,
$\alpha_{\text{GB}} = 0.20$,
$\alpha_{\text{NN}} = 0.15$,
$\alpha_{\text{Stack}} = 0.15$.
We found these through grid search on the WDRA dataset.
To make the score easier for users to understand, we map it
to a percentage:
\begin{equation}
F_{\text{final}}(w) = 0.50 + \left(E(w) \times 0.50\right)
\label{eq:final}
\end{equation}

\noindent
So every recommendation score lands between 50\% and 100\%.
We also compute confidence as:
$C = \max(0.5,\; 1 - 2\sqrt{\sigma^2(\{S_1, \ldots, S_5\})})$.

%%% 4.3.2 LLM Integration %%%

\subsubsection{LLM Integration and Scoring}%
\label{subsubsec:llm-integration}

The language model in our system is
\textbf{Meta Llama 3.3 70B}, which has 70 billion parameters.
We run it with structured JSON output enforcement,
a temperature of 0.3, and a maximum of 4,096 tokens per
response.
The model is given a system prompt that tells it to score
each warehouse according to this formula \cite{b3, b7}:
\begin{equation}
S_{\text{LLM}}(w) = 0.30L + 0.25P + 0.25Z + 0.20Q
\label{eq:llm}
\end{equation}

\noindent
$L$ means location relevance, $P$ is price competitiveness,
$Z$ is size adequacy, and $Q$ is overall quality.
Each one is scored between 0 and 1.

We chose to use only one LLM rather than chaining multiple
models.
Meta Llama 3.3 70B handles all language tasks in SmartSpace.
If the model becomes unavailable, due to network issues or
downtime, the system automatically switches to the local
ML ensemble.
This way, recommendations keep coming without any dependence
on external language model services \cite{b7, b4}.

%%% 4.3.3 Document Verification %%%

\subsubsection{Document Verification System}%
\label{subsubsec:docverify}

When warehouse owners submit documents, the system scores
them automatically:
\begin{equation}
V = 25G + 25P + 15T + 20D + 15C
\label{eq:verify}
\end{equation}

\noindent
$G$ is the GST score, $P$ is PAN verification,
$T$ is phone validation, $D$ is document quality,
and $C$ is completeness.
All are normalized between 0 and 1.
Based on the total, the decision is:
$V \geq 80$ means auto-approve,
$60 \leq V < 80$ means send for review,
$40 \leq V < 60$ means caution,
and $V < 40$ means reject.

%%% 4.4 Comparative Analysis %%%

\subsection{Comparative Analysis: LLM vs ML Ensemble}%
\label{subsec:compare}

We compared the two engines across ten different criteria
to understand where each one does better
\cite{b3, b10}.

\subsubsection{Feature-wise Performance}%
\label{subsubsec:featperf}

Fig.~\ref{fig:feat} shows how each engine performs on six
core platform features.
The ML ensemble does better on tasks where the data is
structured, for example it scores 91\% on recommendation
accuracy and 86\% on search relevance.
The LLM, on the other hand, is stronger at tasks that need
language understanding: it hits 92\% on natural language
booking, 90\% on pricing analysis, and 88\% on document
verification \cite{b7}.

\begin{figure}[h]
\centering
\includegraphics[width=0.95\columnwidth]{Fig1_Feature_ML_vs_LLM_Paper.jpeg}
\begin{center}
    \textbf{Fig. 2} Feature-wise performance comparison between ML and LLM.
\end{center}
\label{fig:feat}
\end{figure}

\subsubsection{Multi-Criteria Evaluation}%
\label{subsubsec:multicrit}

Table~\ref{tab:mlvsllm} gives the full breakdown.
Out of 10 criteria, the LLM scores higher on 8 of them,
with an average of 83.0 compared to 56.6 for the ML ensemble.
But the ML side wins on speed (45 ms versus 850 ms) and cost
(no API fees at all) \cite{b4, b3}.

\begin{table}[h]
\centering
\label{tab:mlvsllm}
\large
\renewcommand{\arraystretch}{1.1}
\begin{tabular}{|l|c|c|}
\hline
\textbf{Criterion} & \textbf{ML} & \textbf{LLM} \\
\hline
Accuracy (\%)        & 91.0 & 88.0 \\
Speed (ms)           & 45   & 850  \\
Scalability          & 75   & 90   \\
NL Understanding     & 15   & 95   \\
Context Awareness    & 20   & 92   \\
Adaptability         & 35   & 88   \\
Maintenance Effort   & 80   & 70   \\
Cost Efficiency      & 95   & 60   \\
Explainability       & 85   & 72   \\
Multi-task           & 25   & 85   \\
\hline
\textbf{Average}     & \textbf{56.6} & \textbf{83.0} \\
\hline
\end{tabular}
\begin{center}
    \textbf{Table 2.} Multi-Criteria Performance Comparison Between ML and LLM.
\end{center}
\end{table}

\begin{figure}[h]
\centering
\includegraphics[width=0.95\columnwidth]{Fig2_Radar_Comparison1.png}
\begin{center}
    \textbf{Fig 3.} Multi-criteria radar comparison across 10 performance dimensions.
LLM achieves 83.0/100 average and ML achieves 56.6/100.
\end{center}
\label{fig:radar}
\end{figure}

%%% 4.5 Model Selection %%%

\subsection{Model Selection: Why Meta Llama 3.3 70B}%
\label{subsec:modelselect}

We did not just pick a language model at random.
We compared Llama 3.3 70B against GPT-4o, Claude 3.5,
Gemini 1.5, and Mistral 7B across six practical parameters.
Table~\ref{tab:models} shows the numbers \cite{b3, b7}.

\begin{table}[h]
\centering
\label{tab:llm-compare}
\renewcommand{\arraystretch}{1.15}
\begin{tabular}{|l|c|c|c|c|c|}
\hline
\textbf{Param.} & \textbf{Llama 3.3} & \textbf{GPT 4o} & \textbf{Claude 3.5} & \textbf{Gemini 1.5} & \textbf{Mistral 7B} \\
\hline
Params (B)      & 70   & 200  & 175  & MoE  & 7   \\
Quality (\%)    & 93   & 96   & 94   & 91   & 78  \\
Speed (t/s)     & 800  & 150  & 180  & 200  & 600 \\
JSON (\%)       & 98   & 95   & 93   & 90   & 82  \\
Avail. (\%)     & 99   & 99.9 & 99.5 & 99   & 95  \\
Latency (ms)    & 850  & 2200 & 1800 & 1500 & 950 \\
Cost/1M tok     & Free & \$5  & \$3  & \$1.25 & Free \\
\hline
\end{tabular}
\begin{center}
    \textbf{Table 3.} LLM Model Comparison Across Operational Parameters.
\end{center}
\end{table}

Three things made Llama 3.3 70B the clear winner for our
use case:
\begin{itemize}
    \item \textbf{Inference Speed:}
      On dedicated inference hardware, Llama runs at about
      800 tokens per second.
      That is 4 to 5 times faster than GPT-4o at 150 tokens
      per second.
      Since our pipeline makes four back-to-back LLM calls,
      this drops the total wait time from 2,200 ms to about
      850 ms, which is a 61\% reduction.

    \item \textbf{JSON Reliability:}
      When we tell the model to return structured JSON,
      it complies 98\% of the time.
      That is the highest among all the models we tested,
      and it matters a lot because our system parses JSON
      responses programmatically.

    \item \textbf{Cost at Scale:}
      With a free inference tier that supports roughly
      14,400 requests per day, we pay nothing for LLM calls.
      GPT-4o would cost about \$2,740 per year for the same
      volume \cite{b3}.
\end{itemize}

\begin{figure}[h]
\centering
\includegraphics[width=0.95\columnwidth]{Fig3_Why_Llama_Model_Comparison.jpeg}
\begin{center}
    \textbf{Fig 4.} Six-panel comparison of Meta Llama 3.3 70B against competing LLMs.
\end{center}
\label{fig:model}
\end{figure}

%%% 4.5.1 Failover %%%

\subsubsection{Failover and Reliability}%
\label{subsubsec:failover}

If the LLM becomes unavailable, whether from a network issue,
a timeout, or any other reason, the system does not go down.
It switches to the local 5-algorithm ML ensemble, which runs
entirely on the server with zero external dependencies.
Since the LLM has a measured uptime of about 99\% and the
local ML fallback is always on, the combined availability
works out to \cite{b4, b7}:
\begin{equation}
A_{\text{sys}}
  = 1 - (1 - A_{\text{LLM}})(1 - A_{\text{ML}})
  = 1 - (0.01)(0)
  = 1.0
\label{eq:avail}
\end{equation}

\noindent
The ML ensemble never goes down because it does not depend
on anything outside the server.
So in practice, SmartSpace always has a working recommendation
engine, regardless of what happens to the language model.

%%% 4.6 Performance Evaluation %%%

\subsection{Performance Evaluation and Results}%
\label{subsec:perf}

We tested the system on the full dataset of
\textbf{10,002 warehouse records} from Maharashtra.
\begin{itemize}
    \item \textbf{Predictive Accuracy:}
      The ML ensemble got a \textbf{mean accuracy of 87.9\%}
      on the test set.
      The \textbf{Precision@10} score was \textbf{0.83},
      which means 8 or 9 out of every top-10 recommendations
      were actually relevant to what the user wanted
      \cite{b1, b10}.

    \item \textbf{Conversational Reliability:}
      Under stress testing, the system handled
      \textbf{96.8\%} of natural language queries
      successfully.
      When we simulated LLM outages, the fallback kept
      response times under 3 seconds
      \cite{b8, b9}.

    \item \textbf{Operational Efficiency:}
      Compared to the old manual approach, SmartSpace
      cut the total search time by about \textbf{60\%}.
      For SMEs and logistics managers who run these
      searches regularly, that is a significant time
      savings \cite{b5, b2}.
\end{itemize}

%% ==========================================================
%% SECTION 5 — CONCLUSION
%% ==========================================================

\section{Conclusion}\label{sec5}

We built and tested SmartSpace, a hybrid ML and LLM-assisted
recommendation platform for warehouses in Maharashtra.
The system works.
Recommendation accuracy went up, search times went down,
and users said the explanations helped them trust the
suggestions.
Compared to the old way of searching through listings
manually, this is a clear step forward
\cite{b2, b5, b8}.

The hybrid approach, with KNN and Random Forest feeding into
a meta-learner, gives us stable rankings that do not flip
around when user preferences change slightly \cite{b10}.
The language model adds a layer of reasoning that pure ML
cannot provide, especially when users describe what they
want in everyday language \cite{b3, b7}.
And the fallback mechanism means the system never goes
offline, even if the LLM has issues \cite{b4}.

There are things we want to improve.
We plan to add online learning so the system adjusts its
weights based on actual user behavior over time.
Better data quality checks are on the list too.
We also want to build a proper benchmarking framework to
measure things more rigorously.
Eventually, we would like to expand beyond Maharashtra,
bring in real-time market data, and add predictive capacity
planning so warehouse operators can plan ahead instead of
reacting to demand after it hits \cite{b1, b8}.

%% ==========================================================
%% ACKNOWLEDGMENT
%% ==========================================================

\section*{Acknowledgment}

We would like to thank the Department of Computer Engineering
at A.P. Shah Institute of Technology, Thane, for providing
the lab facilities and academic guidance we needed for this
project.
We are also grateful to the open-source communities and tool
providers whose frameworks and libraries made it possible
to build and test SmartSpace.

%% ==========================================================
%% REFERENCES — Author names corrected per actual paper PDFs
%% ==========================================================

\begin{thebibliography}{10}

\bibitem{b1}
Z. Xing, Y. Wu, and S. Yu,
``An Enhanced Modelling Approach for Warehouse Sharing Platform
System Designing Problem,''
\textit{arXiv preprint arXiv:2401.15389},
2024.

\bibitem{b2}
M. Matusiewicz and D. Ksi\k{a}\.zkiewicz,
``Shared Logistics---Literature Review,''
\textit{Applied Sciences},
vol.~13, no.~4, p.~2352,
2023.

\bibitem{b3}
T. Wi\'{s}niewski,
``Using large language models (LLMs) to support
simulation-based optimization in supply chain management,''
\textit{Advances in Production Engineering and Management},
vol.~20, no.~4, pp.~491--506,
2025.

\bibitem{b4}
T. Berlec, M. Corn, S. Varljen, and P. Podr\v{z}aj,
``Exploring Decentralized Warehouse Management Using
Large Language Models: A Proof of Concept,''
\textit{Applied Sciences},
vol.~15, no.~10, p.~5734,
2025.

\bibitem{b5}
V. Elia, M. G. Gnoni, and F. Tornese,
``On-Demand Warehousing Platforms: Evolution and Trend Analysis
of an Industrial Sharing Economy Model,''
\textit{Logistics},
vol.~8, no.~4, p.~93,
2024.

\bibitem{b6}
Technical Report,
``Full Stack Systems using React and Supabase,''
2024.

\bibitem{b7}
H. Wang, J. Jiang, L. J. Hong, and G. Jiang,
``LLMs for Supply Chain Management,''
\textit{arXiv preprint arXiv:2505.18597},
2025.

\bibitem{b8}
B. Y. Ekren, U. Venkatadri, F. Sgarbossa,
and E. H. Grosse,
``Warehousing 5.0 for the future of the logistics industry,''
\textit{International Journal of Production Research},
vol.~64, no.~5,
2026.

\bibitem{b9}
C. Liang, L. Sun, S. Luo, X. Liu, and R. Wu,
``Architecture Design of a Smart Logistics Warehouse
Management System based on Multi-modal AI,''
\textit{International Core Journal of Engineering},
vol.~11, no.~5,
2025.

\bibitem{b10}
R. F. de Assis, A. F. Faria,
V. Thomasset-Laperri\`{e}re,
L. A. Santa-Eulalia, M. Ouhimmou,
and W. de~Paula Ferreira,
``Machine Learning in Warehouse Management: A Survey,''
\textit{Procedia Computer Science},
vol.~232, pp.~2790--2799,
2024.

\end{thebibliography}

\end{document}
